\section{Just ask Why?}

In Episode
\#6, titled ``The General''\footnote{\url{http://en.wikipedia.org/wiki/The\_General\_(The\_Prisoner)}} of the \textbf{Prisoner}, there is a scene at an art class where the instructor is teaching creativity. Number 6 asks about the odd activities of three of the students.

\begin{itemize}


\item \emph{One of the students is assiduously tearing out the pages of a book.} The instructor explains that sometimes destruction must precede construction.

\item \emph{One of the students is standing on her head.} That's so she can get a different perspective on things, explains the instructor.

\item \emph{One of the students in napping in his chair.} Well, we can't be creative all the time, is the explanation.
\end{itemize}


Rather than demonstrating techniques to free up creativity, this scene illustrates the principles of social control.

\begin{itemize}


\item Destruction of the social order.

\item Seeing the world opposite to the sane and normal way.

\item Lack of awareness.
\end{itemize}


The episode revolves around a concept called Speedlearn whereby a university level history course is taught in three minutes. It is reminiscent of how ideas and policies that used to be unthinkable not so long ago are now considered normal and desirable. The technique of Speedlearn is well known and in constant use.

In The Prisoner, the process is thwarted. Who can thwart the process now?

Whenever you are told to accept or believe what was never acceptable or believable before, just ask:

\textbf{Why?}


\flrightit{Posted on 2010-04-30 by Cologero}

